\section*{CRediT Authorship Contribution Statement}

Soroush Vahidi: Conceptualization, Methodology, Software, Validation,
Formal analysis, Investigation, Data curation, Writing -- original draft,
Writing -- review and editing, Visualization.

\section*{Declaration of Competing Interest}

The author declares no competing interests.

\section*{Funding}

This work was supported in part by the Google Cloud Research Credits
Program, which provided cloud credits used to support software development
for this study; the CloudRift AI Builder Grant and the Cohere Labs Catalyst
Grant Program, which provided cloud-computing and API credits.

\section*{Acknowledgements}

The author gratefully acknowledges Professor Ioannis Koutis for his research
guidance, continued support, and provision of computational resources that
facilitated this work. Computational resources for this work were provided by
Wulver, the high-performance computing cluster managed by Academic and
Research Computing Systems (ARCS) in the Department of Information Services
and Technology (IST) at the New Jersey Institute of Technology.

\section*{Data Availability}

This study distinguishes four categories of data, consistent with
\S\ref{sec:artifact}. \textbf{Externally sourced public traces} (BurstGPT,
Azure 2024, Bailian/Qwen) are governed by their original publishers' release
terms and license status (\S\ref{sec:provenance}; BurstGPT: CC-BY-4.0) and
are not redistributed as part of this project's release; readers should
obtain them from their original sources. \textbf{Derived benchmark manifests},
\textbf{reproduction scripts}, simulator code, policy implementations, and
analysis code are available in the GitHub repository
\url{https://github.com/SoroushVahidi/llm-serving-scheduler-robustness-benchmark}.
The companion Hugging Face dataset repository
\url{https://huggingface.co/datasets/SoroushVahidi/llm-serving-scheduler-portability}
contains the current public, analysis-ready subset: source-license
documentation, checksums, selected canonical result tables for the principal
portability and sample-complexity results, and the reduced RQ6 analysis
result with computational provenance.
An immutable snapshot of the code repository at the version reported here is
archived with a persistent identifier via Zenodo, DOI
\href{https://doi.org/10.5281/zenodo.22306798}{10.5281/zenodo.22306798}.
The Hugging Face release is an initial, dataset-card-scoped publication
(\S\ref{sec:artifact}); it does \emph{not} include the exploded per-cell
campaign matrix (all $18{,}720$ individual executions), complete per-cell
mechanism telemetry, raw $240$ per-cell RQ6 physical-execution records, or
unpublished intermediate canonical-analysis files consumed by the figure/table
generation script. These non-public records are candidates for a future
dataset revision and are currently represented only by the released reduced
tables, provenance records, and code.

\section*{Code Availability}

The simulator, policy implementations, and analysis pipeline underlying this
study are publicly available at
\url{https://github.com/SoroushVahidi/llm-serving-scheduler-robustness-benchmark},
with an immutable archived snapshot via Zenodo, DOI
\href{https://doi.org/10.5281/zenodo.22306798}{10.5281/zenodo.22306798}.

\section*{Declaration of Generative AI and AI-Assisted Technologies in the Manuscript Preparation Process}

During the preparation of this work, the author used ChatGPT and Codex
(OpenAI), Claude (Anthropic), Gemini (Google), Cursor, and Perplexity AI to
assist with software development, debugging, experimental infrastructure,
literature and research support, and manuscript drafting and editing. After
using these tools, the author reviewed, verified, and revised all AI-assisted
outputs, including code, analyses, references, and manuscript text, as needed
and takes full responsibility for the content of the publication.
